
\section{Introduction}

Transformers~\citep{Vaswani_A_2017_p-nips-OG} have become the dominant architecture in modern deep learning, achieving state-of-the-art results across domains such as natural language processing~\citep{wolf-etal-2020-transformers} and computer vision~\citep{dosovitskiy2021an, zhao2021point, zhai2022scaling}. Despite their success, these models fail on surprisingly simple tasks---such as copying or repeating sequences \citep{Jelassi2024RepeatAM}---once the input exceeds a certain length. These failures persist even for large models specifically trained on those tasks, suggesting that they may reflect intrinsic architectural constraints.

\vspace{1em}
This paper characterizes such constraints by identifying structural limitations that hold for all transformers, independently of model size, training data, or available computation:
\begin{itemize}
    \item We show that every transformer can output only a finite set of output sequences. That is, most sequences are fundamentally inaccessible: no choice of prompt, context length, or inference strategy can cause the model to generate them. This impossibility result provides a principled explanation for well-documented failures of transformers on simple algorithmic tasks, such as exact copying.
    \item Moreover, we show that for a fixed prompt length, the proportion of accessible sequences of length $n$ decays exponentially fast with $n$ beyond a model-dependent threshold. This phenomenon explains a puzzling empirical behavior observed in practice, where transformers perform nearly perfectly up to a critical length and then fail abruptly rather than gradually.
    \item Finally, we derive an explicit formula to predict those thresholds as functions of the transformer architecture being used. Remarkably, this theoretical result closely matches empirical performance observed across a wide range of transformer architectures and model sizes.
\end{itemize}

To establish these results, we analyze how the finite precision and the geometry of the internal representations of a transformer constrain the set of sequences it can represent or generate. We formalize this constraint through the notion of \emph{accessible sequences}, defined as outputs that can arise from some input prompt (Section~\ref{sec:embed}). The main theoretical analysis, which yields the results stated above, is conducted in Section~\ref{sec:main}. We also devise experiments that validate our theoretical predictions in Section~\ref{sec:expe}\footnote{Code is available at \href{https://github.com/mario-michelessa/transformers\_accessibility}{github.com/mario-michelessa/transformers\_accessibility}}.

\begin{remark}[Broader Applicability of Our Results]\label{rmk:mamba}
    Our bounds extend beyond transformers and apply to any architecture with bounded internal representations and finite precision. This may include state-space models such as Mamba~\citep{gu2024mamba}. However, the possibility of diverging hidden states~\citep{lu2026mamba} makes the analysis more challenging, and we leave this direction for future work.
\end{remark}

\subsection{Related Works}

\paragraph{Approximation Theory of Transformers.}
The representational power of transformers has been extensively explored through the lens of approximation theory. In contrast, we study how finite numerical precision constrains this expressivity—showing that, as rounding errors accumulate, even simple functions such as the identity map become difficult to approximate in practice. \citet{Yun_C_2020_p-iclr_continuous} established the first universal approximation property for continuous functions, later extended to sparse attention transformers~\citep{a8b4c6da7b064322905cf50d4103d356} and sparse target functions~\citep{edelman2022inductive}. Subsequent works refined these results for one-layer architectures~\citep{kajitsuka2024are, Jiang_H_2024_p_nips_Jackson-type} and examined approximation under compact output constraints~\citep{kratsios2022universal}. Architectural factors such as relative positional encodings~\citep{luo2022your} have also been analyzed. A complementary line of work investigates approximation properties under prompt tuning~\citep{Wang_Y_2023_p-nips_prompt-og, Hu_J_2025_p-iclr_prompt-smaller}. Among those, \citet{meyer2025memorylimitationsprompttuning} uses a proof strategy similar to ours, leveraging a packing number-type argument.


\paragraph{Empirical Analyses of Sequence Accessibility.}
Empirical studies have also examined the ability of transformers to reproduce or memorize sequences. \citet{barbero2024transformers} showed that transformers struggle to reliably copy long sequences, but their analysis focuses on highly synthetic settings—sequences differing by only a few tokens or alternating between two symbols ({\it e.g.} 0 and 1).
In contrast, our work studies the capacity of transformers to generate, and therefore copy, arbitrary sequences. In this line of work, \citet{Jelassi2024RepeatAM} showed that even transformers trained on the specific task of copying struggle to reliably execute this task as the sequence length increases. The work most closely related to ours is~\citet{kuratov-etal-2025-cramming}, who introduced the “cramming’’ procedure as an exact test of accessibility. For any target sequence $[\mathbf{x}_1,\ldots,\mathbf{x}_n]$, a set of trainable memory vectors $[\mathbf{mem}]=[\mathbf{y}_1,\ldots,\mathbf{y}_m]$ is optimized so that the model generates $[\mathbf{x}_1,\ldots,\mathbf{x}_n]$ when conditioned on $[\mathbf{mem}]$. This approach effectively searches over all possible prompts—including highly unnatural ones—to determine whether a sequence is accessible. Cramming thus provides an operational definition of accessibility: a sequence is accessible if and only if a corresponding encoding $[\mathbf{mem}]$ exists. Our work provides the first rigorous theoretical explanation for the empirical findings in~\citet{kuratov-etal-2025-cramming}, and extends the experimental analysis beyond the original setting. We describe this further in Section~\ref{sec:expe}.

\paragraph{Mean-Field Transformers.}
To analyze prompts of arbitrary length, it is useful to interpret a transformer as a mapping between probability measures rather than finite-length sequences, following the mean-field formulation introduced by~\citet{pmlr-v151-sander22a} and defined in Section~\ref{sec:mean-field}. This perspective has been further developed in several directions. \citet{Castin_V_2024_p-icml_lipschitz} employed it to study the Lipschitz continuity of transformers, while~\citet{furuya2025transformers} extended the approximation theory of transformers to the space of probability distributions. More recently,~\citet{meyer2025memorylimitationsprompttuning} leveraged the mean-field framework to characterize the memory limitations of transformers. Our analysis builds upon this formulation to derive accessibility bounds that hold independent of the prompt.

\paragraph{Representations of the Embedding Space.}
Prior work has visualized and probed LLM embedding spaces in three main ways. First, transformer representations have been shown to be highly anisotropic, with a large fraction of the variance concentrated along a small number of directions~\citep{mu2018all, ethayarajh-2019-contextual}. 
Second, probing methods have shown that certain syntactic and semantic properties are linearly recoverable from transformer hidden representations~\citep{alain2016understanding,hewitt2019structural,kornblith2019similarity}.
Third, decoding intermediate activations by projecting them through the decoder matrix to logits reveals layerwise sharpening of next-token beliefs~\citep{nostalgebraist2020logit,belrose2023eliciting}.
Inspired by this decoding view, we use the same decoder matrix to find hyperplanes that partition the embedding space into next-token argmax cells in Section~\ref{sec:embed}.

\paragraph{Expressive Power of Transformers.}
Complementary to the line of work on approximation theory, the expressivity of transformers can also be studied in terms of the formal languages they can express~\citep{strobl-etal-2024-formal}. Transformers have been shown to represent Turing machines~\citep{NEURIPS2022_4ebf1d74} and arbitrary computer programs~\citep{pmlr-v202-giannou23a}. The impact of chain-of-thought reasoning~\citep{merrill2024the} and precision~\citep{chiang2025transformers} have also been studied. Most relevant to our work is the work of~\citet{huang2025a}, who prove the impossibility of copying arbitrary sequences for transformers with absolute positional encodings and under some idealized assumptions. We extend those results by removing those assumptions and providing both a qualitative and quantitative analysis of the performance of transformers in the copying task.

\subsection{Notations}

Bold lowercase letters (e.g., $\mathbf{x}$) denote vectors, and bold uppercase letters (e.g., $\mathbf{W}$) denote matrices. For a matrix $\mathbf{W}$, we write $\mathbf{W}_{i,j}$, $\mathbf{W}_{i,:}$, and $\mathbf{W}_{:,j}$ to refer to its $(i,j)$-th element, $i$-th row, and $j$-th column, respectively. We can use $i=-1$ and $j=-1$ to denote the last row and column, respectively. The concatenation of $k$ matrices with the same number of  rows  $d$, $\mathbf W^1\in\mathbb R^{d\times m_i},\dots,\mathbf W^k\in\mathbb R^{d\times m_k}$ is denoted as $[\mathbf W^1,\dots,\mathbf W^k]\in\mathbb R^{d\times\sum_{i=1}^km_i}$.

We write $\sigma$ to denote the softmax function. The rectified linear unit is defined as $\mathrm{ReLU}(v) = \max(v, 0)$, where the maximum is applied elementwise. We use $|\mathcal A|$ to denote both the cardinality of a finite set, and the volume of a set $\mathcal A\subset\mathbb R^n$ for some $n\in\mathbb N$ (we take $\mathbb N$ to be the set of strictly positive integers). We denote by $\mathcal V=\{t_1,\ldots,t_{|\mathcal V|}\}$ the vocabulary (that is the set of all distinct tokens) and $d$ the embedding dimension. The probability simplex over the vocabulary is denoted as $\Delta^{\mathcal{V}}$.

Throughout, $\|\cdot\|$ denotes a norm, which by default may be taken as the standard $\ell_\infty$ or $\ell_p$ norm for vectors. Unless otherwise specified, our results hold for all of these norms. Let $\Pi(\mu, \nu) $ be the set of all joint distributions with marginals $\mu $ and $\nu$; this is also called the set of {\em couplings} of $\mu$ and $\nu$. When considering distributions, we use the Wasserstein distance $$W_q(\mu, \nu) = \left( \inf_{\pi \in \Pi(\mu, \nu)} \int \|x - y\|^q \, \mathrm{d}\pi(x, y) \right)^{1/q}.$$ This distance is further discussed in Appendix~\ref{app:wasserstein}.